%%%
%
% $Autor: Wings $
% $Datum: 2021-05-14 $
% $Pfad: GitLab/MLEdgeComputer $
% $Dateiname: ArduinoIDESetupAlvik
% $Version: 4620 $
%
% !TeX spellcheck = de_DE/GB
% !TeX program = pdflatex
% !BIB program = biber/bibtex
% !TeX encoding = utf8
%
%%%



\chapter{Arduino IDE 2 Setup für den Arduino Alvik} \Mynote{Deutsch}


\section{Einführung}

Arduino IDE 2 ist die offizielle Entwicklungsumgebung von Arduino zum Schreiben, Kompilieren und Hochladen von Code auf Mikrocontroller-Boards. Sie bietet Funktionen wie Code-Hervorhebung, Autovervollständigung und Fehlerwarnungen. Zusätzlich enthält sie integrierte Tools für Debugging, den seriellen Monitor sowie die Verwaltung von Boards und Bibliotheken. Die IDE ist mit offiziellen Arduino-Boards sowie ausgewählten Drittanbieter-Boards kompatibel.

\section{Installation}

\begin{enumerate}
	\item\textbf{Einrichten Arduino IDE} 
	\begin{itemize}
		
		\item Installieren die Arduino IDE von der offiziellen Arduino-Website.
		\item Öffnen die Arduino IDE.Gehen zu \menu{Sketch > Include Library > Manage Libraries}.
		\item Suchen im Bibliotheksmanager nach ``Alvik'' und installieren die neueste Version der Alvik-Bibliothek.
	\end{itemize} 
	
	
	\item\textbf{Alvik für die Arduino IDE vorbereiten}
	\begin{itemize}
		
		\item Verbinden Pin B1 mit GND auf der Alvik Board.
		\item Während beide Pins verbunden sind, schlie\ss en das Alvik-Board über das USB-Kabel an den Computer an.
		\item Suchen im Bibliotheksmanager nach „Alvik“ und installieren die neueste Version der Alvik-Bibliothek.
		\item Wählen esptool als Programmer im Menü \menu{Tools > Programmer} aus. Das Verfahren ist in der Abbildung \ref{EsptoolSelection} dargestellt.
		\begin{figure}[ht]  
			\begin{center}
				\includegraphics[width=12cm]{Arduino/Alvik/EsptoolSelection}
				\caption{Esptool-Auswahl} 
				\label{EsptoolSelection}                                                                                                                                                                                                                                                                                                          
			\end{center}
		\end{figure}
		
		\item Wählen Upload Using Programmer aus dem Sketch-Menü.Das Verfahren ist in der Abbildung\ref{UploadAlvikProgrammer} dargestellt.Drücken die Reset-Taste auf der Platine, um sicherzustellen, dass sie zum Hochladen bereit ist.
		\begin{figure}[ht]  
			\begin{center}
				\includegraphics[width=12cm]{Arduino/Alvik/UploadAlvikProgrammer.png}
				\caption{Hochladen mit Programmierer-Option} 
				\label{UploadAlvikProgrammer}                                                                                                                                                                                                                                                                                                          
			\end{center}
		\end{figure}
		
		\Mynote{Beschreiben der Bilder, Beispiel Code als Code einbinden, Bibliothek beschreiben + Funktionen nach Template }
		\item Öffnen das Drive-Beispiel in der Arduino IDE, indem zu File > Examples > Arduino\_Alvik > drive gehen.Testet die Alvik-Verbindung und bringt den Alvik-Motor in Gang.Das Verfahren ist in der Abbildung\ref{uploadExample} dargestellt.
		\begin{figure}[ht]  
			\begin{center}
				\includegraphics[width=12cm]{Arduino/Alvik/uploadExample.png}
				\caption{Hochladen von Testprogrammen} 
				\label{uploadExample}                                                                                                                                                                                                                                                                                                          
			\end{center}
		\end{figure}
	\end{itemize}
	
\end{enumerate}


\section{Erste Schritte}

Nachdem Alvik korrekt eingerichtet ist, können einige Sketch hochladen.
\begin{enumerate}
	\item\textbf{Bibliotheken} 
	\begin{itemize}
		\item\textbf{Arduino\_Alvik:} Dies ist die Hauptbibliothek, die im Sketch verwendet wird und Befehle zur Steuerung des Arduino Nano ESP32-Boards in Alvik enthält.
		\item\textbf{Arduino\_AlvikCarrier:} Diese Bibliothek wurde für das STM-Board auf dem Gerät entwickelt und ist in Situationen nützlich, in denen eine feinere Steuerung der Befehle erforderlich ist. Sie ermöglicht komplexere Entwicklungen, insbesondere wenn eine tiefere Integration mit der Hardware erforderlich ist.
	\end{itemize} 
	\item\textbf{Print Firmware Version:} 
	\begin{itemize}
		
		\item Die Arduino\_Alvik Bibliothek muss zuerst eingebunden sein. Diese Bibliothek enthält alle vordefinierten Funktionen und Klassen, die zur Vereinfachung der Steuerung der Roboterhardware, wie Motoren und Sensoren, benötigt werden. Ohne die include-Anweisung erkennt der Compiler die Klasse Arduino\_Alvik und die zugehörigen Methoden nicht.Der Code ist in der Abbildung\ref{BibliothekCode} dargestellt.
		\begin{figure}[ht]  
			\begin{center}
				\includegraphics[width=12cm]{Arduino/Alvik/BibliothekCode.png}
				\caption{Arduino\_Alvik Bibliothek Code} 
				\label{BibliothekCode}                                                                                                                                                                                                                                                                                                          
			\end{center}
		\end{figure}
		\item Um das Alvik-Livarium zu verwenden, deklarieren Sie ein Objekt der Klasse Arduino\_Alvik, in diesem Fall mit dem Namen alvik. Alle Interaktionen mit dem Roboter finden über dieses alvik-Objekt statt.Der Code ist in der Abbildung\ref{alvikKlassencode} dargestellt.
		\begin{figure}[ht]  
			\begin{center}
				\includegraphics[width=12cm]{Arduino/Alvik/alvikKlassencode.png}
				\caption{alvik Klassen code} 
				\label{alvikKlassencode}                                                                                                                                                                                                                                                                                                          
			\end{center}
		\end{figure}
		\item In der Funktion setup() initialisiert die Methode alvik.begin() den Alvik-Roboter. Dies ist bei allen Sketch für den Alvik erforderlich.Serial.begin(115200) ist der Initialisierungsschritt für die serielle Kommunikation, der die Baudrate auf 115200bps setzt.Der Code ist in der Abbildung\ref{Setupcode} dargestellt.
		\begin{figure}[ht]  
			\begin{center}
				\includegraphics[width=12cm]{Arduino/Alvik/Setupcode.png}
				\caption{setup code} 
				\label{Setupcode}                                                                                                                                                                                                                                                                                                          
			\end{center}
		\end{figure}
		\item  Mit Hilfe des Codes in der Abbildung\ref{AlvikTestCode} wird die Firmware-Version kontinuierlich abgefragt und jede Sekunde über get\_version() ausgegeben, während die Räder des Alvik-Roboters mit dem Befehl drive() hin und her gedreht werden.
		\begin{figure}[ht]  
			\begin{center}
				\includegraphics[width=12cm]{Arduino/Alvik/AlvikTestCode.png}
				\caption{Alvik Test Code} 
				\label{AlvikTestCode}                                                                                                                                                                                                                                                                                                          
			\end{center}
		\end{figure}
	\end{itemize}
	
	
\end{enumerate}

